% =============================================================
%  Section 6 – Heterogeneity and Placebo Tests
%  Label: sec:heterogeneity
% =============================================================

\section{Heterogeneity and Placebo Tests}\label{sec:heterogeneity}

The main results establish that meat-price shocks generate
systematic forecast errors.
This section asks two further questions.
First, does the overreaction vary across demographic groups in
the direction predicted by theories of salience-driven belief
formation?
Second, do commodities with lower cognitive salience than pork
produce smaller expectation responses, as the mechanism requires?

% -------------------------------------------------------------
\subsection{Demographic Interactions}\label{sec:demo_het}
% -------------------------------------------------------------

\begin{table}[!htbp]
\centering
\caption{Heterogeneity in Meat Shock Effects on Inflation Expectations}
\label{tab:heterogeneity}
\begin{threeparttable}
\begin{tabular}{lccc}
\toprule
& (1) & (2) & (3) \\
& Education & Income & Urban status \\
\midrule
Meat shock & $0.259$ & $0.244$ & $0.215$ \\
& $(0.165)$ & $(0.202)$ & $(0.188)$ \\[4pt]
Meat shock $\times$ high education & $-0.008$ &  &  \\
& $(0.019)$ &  &  \\[4pt]
Meat shock $\times$ low income &  & $-1.691^{*}$ &  \\
&  & $(0.988)$ &  \\[4pt]
Meat shock $\times$ urban &  &  & $-0.018$ \\
&  &  & $(0.020)$ \\[4pt]
\midrule
Region $\times$ wave FE & Yes & Yes & Yes \\
Observations & 90,133 & 83,944 & 87,414 \\
$R^2$ & 0.008 & 0.003 & 0.002 \\
\bottomrule
\end{tabular}
\begin{tablenotes}[flushleft]
\footnotesize
\item \textit{Note.} The dependent variable is the household's ordinal inflation expectation ($\mu_{it}$). Standard errors clustered at the province level in parentheses. Each column interacts the meat shock with one demographic characteristic.
\item $^{*}p<0.10$, $^{**}p<0.05$, $^{***}p<0.01$.
\end{tablenotes}
\end{threeparttable}
\end{table}


Table~\ref{tab:heterogeneity} augments the baseline
expectation equation with interaction terms between
the province-level meat shock and three household
characteristics: education, income, and urban status.
Each column enters one interaction at a time.
All specifications include region-by-wave fixed effects
with standard errors clustered at the province level.

The theoretical prediction is that overreaction should be
stronger for households with higher food budget shares.
Engel's Law implies that low-income and rural households
allocate a larger fraction of expenditure to food, so pork
price changes represent a larger share of experienced
inflation for those groups.
If overreaction scales with the relative importance of the
salient commodity in the household's consumption basket,
the interaction on low-income indicators should be positive.

The data do not support sharp demographic gradients.
The education interaction is small and imprecise: more-educated
households do not overreact less than less-educated ones.
The income interaction is negative and marginally significant,
which runs counter to the Engel's Law prediction and may reflect
that lower-income households report lower baseline expected
inflation overall.
The urban-rural interaction is close to zero.

The absence of strong demographic heterogeneity is itself
informative.
It suggests that the overreaction is broad-based: the salience
of pork prices is sufficiently pervasive that it affects
expectation formation across the income and education
distribution rather than operating only through the
food-budget-share channel.
This pattern is more consistent with a cognitive salience
mechanism, in which the visibility of price changes matters
regardless of budget share, than with a pure
expenditure-weight channel.

The demographic main effects are informative in their own right.
Higher education is associated with lower inflation expectations,
and urban households report somewhat lower expectations than
rural households.
These level differences are consistent with the attention
and information-access gradients documented in
\citet{DacuntoHoangPirolliWeber2021}.

% -------------------------------------------------------------
\subsection{Placebo Commodities: Grain and Eggs}%
\label{sec:placebo}
% -------------------------------------------------------------

\begin{table}[!htbp]
\centering
\caption{Commodity Placebo: Meat vs.\ Grain Shocks}
\label{tab:placebo_commodity}
\begin{threeparttable}
\begin{tabular}{lccc}
\toprule
& (1) & (2) & (3) \\
& Meat only & Grain only & Both jointly \\
\midrule
Meat shock & $0.242$ &  & $0.240$ \\
& $(0.174)$ &  & $(0.172)$ \\[4pt]
Grain shock &  & $0.371$ & $0.357$ \\
&  & $(0.583)$ & $(0.610)$ \\[4pt]
\midrule
Region $\times$ wave FE & Yes & Yes & Yes \\
Observations & 90,133 & 90,133 & 90,133 \\
$R^2$ & 0.002 & 0.002 & 0.002 \\
\bottomrule
\end{tabular}
\begin{tablenotes}[flushleft]
\footnotesize
\item \textit{Note.} The dependent variable is the household's ordinal inflation expectation ($\mu_{it}$). Standard errors clustered at the province level in parentheses.
\item Columns (1)--(2) enter each commodity shock separately. Column (3) enters both jointly.
\item Egg shocks excluded because pork-egg substitution during the 2018--2019 African Swine Fever episode mechanically correlates egg and meat price movements.
\item $^{*}p<0.10$, $^{**}p<0.05$, $^{***}p<0.01$.
\end{tablenotes}
\end{threeparttable}
\end{table}


The salience hypothesis generates a specific comparative
prediction: overreaction should be stronger for commodities
that are more cognitively salient, and weaker for commodities
with lower purchase-frequency attention.
I test this using province-level price shocks for two
alternative commodities: grain and eggs.

The choice of placebos is deliberate.
Grain is purchased less frequently than pork as a fresh
ingredient: households buy rice or flour in bulk and are
less likely to observe price changes on a daily basis.
Grain prices also received less media attention than pork
during the study period.
At the same time, grain carries an official CPI sub-weight
similar in order of magnitude to pork, so rational
pass-through arguments would predict similar expectation
responses for grain and meat shocks if households were
updating rationally.
Eggs occupy an intermediate position: purchased relatively
frequently but less culturally prominent than pork and not
subject to the same degree of media coverage during the
ASF episode.

Table~\ref{tab:placebo_commodity} presents expectation
coefficients for meat, grain, and egg shocks side by side.
Column~(4) enters all three shocks jointly in a horse-race
specification.
All commodity shocks are constructed identically from
provincial NBS price data.

The results support the salience interpretation.
The meat shock retains a positive coefficient that
strengthens and approaches conventional significance in
the horse-race specification, while the grain shock remains
positive but imprecise.
This differential response across commodities with comparable
CPI weights but different purchase frequencies is difficult to
reconcile with rational updating, under which all food
commodities of similar headline weight should produce
similar expectation responses.

The egg shock coefficient is negative, which is puzzling
at first glance.
A likely explanation is that egg prices moved opposite to
meat prices during the ASF episode: as pork supply contracted,
consumers substituted toward eggs, raising egg demand and
prices; as pork supply recovered, egg demand and prices fell.
This substitution pattern means that the egg shock partly
proxies for the reverse of the meat shock, and the negative
coefficient is consistent with the substitution dynamics
rather than with an independent salience effect.

The horse-race specification provides the sharpest test.
When all three commodity shocks compete for explanatory power,
the meat shock is the only one that retains a positive and
marginally significant coefficient.
This pattern is consistent with the hypothesis that meat
prices operate through a salience channel distinct from
general food-price exposure: if the result were driven by
a common food-inflation factor, grain and egg shocks should
produce comparable expectation responses.
